\documentclass[11pt]{article}
\usepackage[a4paper,margin=27mm]{geometry}
\usepackage{amsmath,amssymb,amsthm,hyperref}
\newtheorem{theorem}{Theorem}
\newtheorem{lemma}[theorem]{Lemma}
\newtheorem{corollary}[theorem]{Corollary}
\newcommand{\Lie}{\operatorname{Lie}}
\newcommand{\sgn}{\operatorname{sgn}}
\title{The exact support barrier for universal relabeling averages}
\author{An application of Klyachko's classical constituent theorem}
\date{30 September 2026}
\begin{document}
\maketitle
Let $L_{n,k}$ be the degree-$k$ part of the reduced free Lie algebra on
$n$ labeled generators, where every bracket with a repeated letter
vanishes. The symmetric group $S_n$ acts by relabeling generators.
For a probability measure $\mu$ on $S_n$, put
\[
 T_\mu=\sum_{g\in S_n}\mu(g)g.
\]
Call $\mu$ a universal relabeling annihilator if $T_\mu$ acts as zero on
$L_{n,k}$ for every $2\le k\le n$.

\begin{theorem}\label{main}
For $n\ge4$, the universal relabeling annihilators are exactly
\[
 \mu(g)=\frac{1+c\,\sgn(g)}{n!}\qquad(-1\le c\le1).
\]
Their minimum support size is $n!/2$, attained only by the uniform
measure on the even permutations or on the odd permutations.
For $n=2,3$, the unique annihilator is the uniform measure on $S_n$.
For $n\ge7$, and also for $n=5$, annihilating just the top degree
$L_{n,n}$ already forces the displayed formula.
\end{theorem}

\begin{lemma}[Constituents]\label{constituents}
For $n\ge4$, the union of the irreducible constituents of the modules
$L_{n,k}$, $2\le k\le n$, consists of every irreducible $S_n$ module
except the trivial and sign representations.
For $n=2,3$, it consists of all nontrivial irreducibles.
\end{lemma}
\begin{proof}
Write $\Lie_k$ for the multilinear free Lie representation of $S_k$.
Decomposition by the set of letters occurring in a bracket gives
\[
 L_{n,k}\cong
 \operatorname{Ind}_{S_k\times S_{n-k}}^{S_n}
       (\Lie_k\boxtimes\mathbf1).
\]
Klyachko's constituent theorem says that for $k>1$, $\Lie_k$ contains
all Specht modules of size $k$ except the trivial module, the sign
module when $k>2$, and the two exceptional shapes $(2,2)$ and
$(2,2,2)$. See \cite[Theorem 1.1, Theorem 1.2 and Remark 1.3]{swanson}
for the exact classification and its connection to the multilinear
Lie representation, and \cite{kovacs} for a combinatorial proof of
the Lie constituent theorem.
Thus for $n\ge7$ or $n=5$, the top degree alone contains every
irreducible except trivial and sign. For $n=4$, the missing shape
$(2,2)$ occurs in
$L_{4,3}=\operatorname{Ind}_{S_3}^{S_4}\Lie_3$, since
$\Lie_3=S^{(2,1)}$ and one-box branching adds $(2,2)$.
For $n=6$, the missing shape $(2,2,2)$ occurs in $L_{6,5}$ since
$S^{(2,2,1)}$ occurs in $\Lie_5$ and one-box branching applies.

No $L_{n,k}$ contains the trivial representation: Frobenius reciprocity
reduces its multiplicity to that of the trivial representation in
$\Lie_k$, which is zero for $k\ge2$.
For $n\ge4$, none contains sign either. If $n-k\ge2$, its restriction
has nontrivial sign on the $S_{n-k}$ factor, incompatible with
$\mathbf1$ there. If $n-k\le1$, then $k\ge3$, and $\Lie_k$ itself
has no sign constituent. This proves the claim for $n\ge4$.
For $n=2$, $\Lie_2$ is sign. For $n=3$, $L_{3,3}$ is the standard
module, while $L_{3,2}=\operatorname{Ind}_{S_2}^{S_3}\sgn$ contains
both standard and sign.
\end{proof}

\begin{proof}[Proof of Theorem~\ref{main}]
Work over $\mathbb C$. If an irreducible module occurs in a module on
which $T_\mu$ vanishes, the matrix Fourier transform
$\sum_g\mu(g)\rho(g)$ is zero on that irreducible. Conversely,
vanishing on every constituent implies vanishing on the whole module.
Lemma~\ref{constituents} therefore says that for $n\ge4$, all Fourier
components of $\mu$ vanish except possibly those belonging to trivial
and sign. Fourier inversion on the finite group yields
\[
 \mu(g)=\frac{1+c\,\sgn(g)}{n!},
 \qquad c=\sum_g\mu(g)\sgn(g).
\]
Normalization supplies the coefficient one of the trivial component;
nonnegativity is equivalent to $-1\le c\le1$. For $|c|<1$ all $n!$
values are positive; for $c=\pm1$ exactly one parity class is supported.
Every such measure annihilates the listed constituents, proving the
converse and the support assertion. The small cases and top-degree
assertion follow from the same constituent calculation.
\end{proof}

\paragraph{Consequences and limits.}
There is no universal positive relabeling average with support
$\exp(O(n))$. In fact its support is at least $n!/2$, even if only the
highest-degree error must be annihilated when $n\ge7$.
This is a barrier for an operator required to annihilate \emph{every}
possible error in the stated Lie modules. It is not a lower bound for
an average tailored to a single error vector, for nonlinear operations,
or for the size of permutation-fair dice.

Uniform averaging over $A_n$ can replace averaging over all $S_n$ in
the universal fair-prefix completion construction when $n\ge4$.
The identity permutation is present, so the literal prefix is preserved.
This halves each symmetrization multiplier but does not change its
factorial asymptotic scale. Relabeling generators is the group action
used throughout; rotation of positions inside a word is a different
operation and is not covered by this theorem.

\paragraph{Exact finite checks.}
The script \texttt{check\_lie\_constituents.py} independently counts
standard tableaux by major index using corner recursion. For every
$2\le n\le20$, it verifies the stated exceptional shapes and the
identity $\sum_\lambda d_\lambda m_\lambda=(n-1)!$.
These computations corroborate the cited theorem; the proof above
does not depend on a finite-range extrapolation.

\paragraph{Attribution.}
The constituent theorem and finite-group Fourier decomposition are
classical. The statement above is their direct application to the
universal averaging mechanism in the dice problem. Historical novelty
of this precise application has not been established.

\begin{thebibliography}{9}
\bibitem{swanson} Joshua P.~Swanson, \emph{On the existence of tableaux
with given modular major index}, Algebraic Combinatorics 1 (2018), 3--21,
\href{https://doi.org/10.5802/alco.4}{doi:10.5802/alco.4}.
\bibitem{kovacs} L.~G.~Kov\'acs and Ralph St\"ohr,
\emph{A combinatorial proof of Klyachko's theorem on Lie representations},
Journal of Algebraic Combinatorics 23 (2006), 225--230,
\href{https://doi.org/10.1007/s10801-006-7394-6}{doi:10.1007/s10801-006-7394-6};
\href{https://archives.maths.anu.edu.au/people/Kovacs/K111.pdf}{author's full text}.
\end{thebibliography}
\end{document}
